\documentclass[a4paper]{article}

\usepackage{hyperref}
\usepackage[utf8]{inputenc}
\usepackage[a4paper, margin=3.5cm]{geometry}
\usepackage[spanish,es-tabla]{babel}
\usepackage{graphicx}
\usepackage{url}
\usepackage{fancyhdr}
\usepackage{float}
\usepackage[newfloat]{minted}
\usepackage[
    type={CC},
    modifier={by-nc-sa},
    version={3.0},
]{doclicense}

\renewcommand{\headrulewidth}{0.6pt}
\renewcommand{\footrulewidth}{0.6pt}

\pagestyle{fancy}
\setlength\headheight{50pt}
\lhead{\includegraphics[height=1.5cm]{logos/upm_logo}}
\chead{Bases de Datos\\\vspace{.5em} Soluciones a ejercicios SQL\\\vspace{-.1em}}
\rhead{\includegraphics[height=1.5cm]{logos/etsisi_logo}}
\lfoot{\textbf{Tema 4:} El lenguaje SQL}
\cfoot{}
\rfoot{\thepage}

\parskip 1.1ex % paragraph spacing

\begin{document}

\section*{Carreras de Formula 1}

Soluciones a las consultas SQL:

\begin{enumerate}
    \item % Obtener el nombre de los pilotos que han participado en al menos 1 carrera del año 2016.
    \begin{minted}[fontsize=\scriptsize, frame=lines, xleftmargin=\parindent, breaklines, framesep=2mm]{SQL}  
SELECT DISTINCT drivers.surname, drivers.forename
FROM drivers 
	INNER JOIN results ON drivers.driverId = results.driverId
	INNER JOIN races ON results.raceId = races.raceId
WHERE year = 2016
    \end{minted}
    
    \item % Nombre de los constructores con los que han disputado carreras más de 50 pilotos diferentes.
    \begin{minted}[fontsize=\scriptsize, frame=lines, xleftmargin=\parindent, breaklines, framesep=2mm]{SQL}  
SELECT name
FROM constructors
WHERE constructorId IN (SELECT constructorId
                        FROM results
                        GROUP BY constructorId
                        HAVING COUNT(DISTINCT driverId) > 50)
    \end{minted}
    
    \item % Nombre y apellidos de los pilotos que nunca han ganado una carrera.
    \begin{minted}[fontsize=\scriptsize, frame=lines, xleftmargin=\parindent, breaklines, framesep=2mm]{SQL}  
SELECT forename, surname
FROM drivers
WHERE driverId NOT IN (SELECT driverId
               FROM results
               WHERE positionOrder = 1)
    \end{minted}            
    
    \item % Obtener el nombre y apellidos de los pilotos que durante el año 2017 han participado en todas las carreras.
    \begin{minted}[fontsize=\scriptsize, frame=lines, xleftmargin=\parindent, breaklines, framesep=2mm]{SQL}  
-- Solución con HAVING COUNT
SELECT drivers.surname, drivers.forename
FROM drivers
WHERE driverId IN (SELECT driverId
                   FROM results
                   WHERE raceId IN (SELECT raceId FROM races WHERE year=2017)
                   GROUP BY results.driverId
                   HAVING COUNT(distinct results.raceId) = (SELECT COUNT(*)
                                                            FROM races
                                                            WHERE year=2017))
                                
-- Solución con NOT EXISTS                                     
SELECT drivers.surname, drivers.forename
FROM drivers
WHERE NOT EXISTS (SELECT *
            FROM races
            WHERE year = 2017 
              AND NOT EXISTS (SELECT *
                              FROM results
                              WHERE results.driverId = drivers.driverId 
                                AND races.raceId = results.raceId))
    \end{minted}
    
    \item % Obtener el nombre, localización, país y año para cada circuito de las carreras que se han disputado entre 2015 y 2017, ordenado por el id del circuito.
    \begin{minted}[fontsize=\scriptsize, frame=lines, xleftmargin=\parindent, breaklines, framesep=2mm]{SQL}  
SELECT circuits.circuitId, circuits.name, circuits.location, circuits.country, year
FROM races 
    INNER JOIN circuits ON races.circuitId = circuits.circuitId
WHERE year BETWEEN 2015 AND 2017
ORDER BY circuits.circuitId
    \end{minted}
    
    \item % Obtener los constructores que no han participado en alguna clasificación.
    \begin{minted}[fontsize=\scriptsize, frame=lines, xleftmargin=\parindent, breaklines, framesep=2mm]{SQL}           
SELECT *
FROM constructors
WHERE constructors.constructorId NOT IN (SELECT qualifying.constructorId 
                                         FROM qualifying)  
    \end{minted}            
    
    \item % Obtener nombres, apellidos de los pilotos que han ganado más de 30 grandes premios así como el número de grandes premios que han ganado.
    \begin{minted}[fontsize=\scriptsize, frame=lines, xleftmargin=\parindent, breaklines, framesep=2mm]{SQL}  
SELECT drivers.surname, drivers.forename , tabla1.Ganado
FROM (SELECT .driverId , COUNT(*) Ganado
FROM `results`
WHERE positionOrder = 1
GROUP BY driverId
HAVING COUNT(*) > 30 ) tabla1 JOIN drivers ON tabla1.driverId=drivers.driverId 
    \end{minted}    
    
    \item % Nombre y apellidos del piloto que obtuvo la vuelta con velocidad media más alta así como el circuito y el año en el que se obtuvo.
    \begin{minted}[fontsize=\scriptsize, frame=lines, xleftmargin=\parindent, breaklines, framesep=2mm]{SQL}  
SELECT drivers.forename, drivers.surname, circuits.name, races.year, results.fastestLapSpeed
FROM drivers
    INNER JOIN results ON drivers.driverId = results.driverId
    INNER JOIN races ON races.raceId = results.raceId
    INNER JOIN circuits ON circuits.circuitId = races.circuitId
WHERE fastestLapSpeed >= ALL (SELECT fastestLapSpeed
                              FROM results)
    \end{minted}
    
    \item % Obtener el nombre, apellidos y la velocidad media del piloto que obtuvo la vuelta con velocidad media más alta en el gran premio de Japón de 2009.
    \begin{minted}[fontsize=\scriptsize, frame=lines, xleftmargin=\parindent, breaklines, framesep=2mm]{SQL}                
SELECT drivers.forename, drivers.surname, results.fastestLapSpeed
FROM drivers
	INNER JOIN results ON results.driverId = drivers.driverId
    INNER JOIN races ON races.raceId = results.raceId
WHERE races.year = 2009
  AND races.name = 'Japanese Grand Prix'
  AND results.fastestLapSpeed >= ALL (SELECT fastestLapSpeed
									  FROM results
										INNER JOIN races ON races.raceId = results.raceId
									  WHERE year = 2009
                                        AND name = 'Japanese Grand Prix')
    \end{minted}             

    \item % Obtener el nombre de los pilotos que durante el año 2017 consiguieron puntos en todas las carreras.
    \begin{minted}[fontsize=\scriptsize, frame=lines, xleftmargin=\parindent, breaklines, framesep=2mm]{SQL}  
-- Solución con HAVING COUNT        
SELECT drivers.forename, drivers.surname
FROM results
	INNER JOIN races ON races.raceId = results.raceId
    INNER JOIN drivers ON drivers.driverId = results.driverId
WHERE points > 0
  AND year = 2017
GROUP BY drivers.driverId, drivers.forename, drivers.surname
HAVING COUNT(*) = (SELECT COUNT(*)
				   FROM races
                   WHERE year = 2017)

                             
-- Solución con NOT EXISTS 
SELECT drivers.surname, drivers.forename
FROM drivers
WHERE NOT EXISTS ( SELECT *
            FROM races
            WHERE year = 2017 AND NOT EXISTS (SELECT *
                                              FROM results
                                              WHERE results.points > 0 
                                                AND drivers.driverId = results.driverId 
                                                AND results.raceId = races.raceId ))
    \end{minted}  
    
    \item % Obtener el nombre de los pilotos, el circuito y el total de paradas, para aquellos pilotos que entre todos los grandes premios disputados han realizado en alguno de ellos el mayor número de paradas y también los que han realizado el menor número de ellas.
    \begin{minted}[fontsize=\scriptsize, frame=lines, xleftmargin=\parindent, breaklines, framesep=2mm]{SQL}   
-- Solución con GROUP BY
SELECT drivers.forename, drivers.surname, circuits.name, races.year, COUNT(*)
FROM drivers
	INNER JOIN pitstops ON pitstops.driverId = drivers.driverId
    INNER JOIN races ON races.raceId = pitstops.raceId
    INNER JOIN circuits ON circuits.circuitId = races.circuitId
GROUP BY drivers.driverId, drivers.forename, drivers.surname, races.raceId, races.year, circuits.name
HAVING COUNT(*) >= ALL(SELECT COUNT(*)
					   FROM pitstops
					   GROUP BY driverId, raceId)

UNION

SELECT drivers.forename, drivers.surname, circuits.name, races.year, COUNT(*)
FROM drivers
	INNER JOIN pitstops ON pitstops.driverId = drivers.driverId
    INNER JOIN races ON races.raceId = pitstops.raceId
    INNER JOIN circuits ON circuits.circuitId = races.circuitId
GROUP BY drivers.driverId, drivers.forename, drivers.surname, races.raceId, races.year, circuits.name
HAVING COUNT(*) <= ALL(SELECT COUNT(*)
					   FROM pitstops
					   GROUP BY driverId, raceId)

-- Solución con el atributo stop
SELECT d.forename, d.surname, c.name, r.year, p.stop
FROM drivers d
	INNER JOIN pitstops p ON p.driverId = d.driverId
    INNER JOIN races r ON r.raceId = p.raceId
    INNER JOIN circuits c ON c.circuitId = r.circuitId
WHERE p.stop = (SELECT MAX(stop)
				FROM pitstops)
   OR p.stop = (SELECT MIN(stop)
				FROM pitstops)
    \end{minted}            
    
    \item % De entre todos los pilotos que han participado en todas las rondas de clasificación (Q1, Q2 y Q3) del gran premio de Abu Dhabi de 2017 (\texttt{qualifying.q1 <> '' AND qualifying.q2 <> '' AND qualifying.q3 <> ''}), obtener el nombre de los pilotos y el id de los equipos, para aquellos equipos que tienen a sus dos pilotos en esa situación.
    \begin{minted}[fontsize=\scriptsize, frame=lines, xleftmargin=\parindent, breaklines, framesep=2mm]{SQL}             
SELECT drivers.forename, drivers.surname, constructors.name
FROM drivers
	INNER JOIN qualifying ON qualifying.driverId = drivers.driverId
    INNER JOIN constructors ON qualifying.constructorId = constructors.constructorId
    INNER JOIN races ON races.raceId = qualifying.raceId
WHERE races.year = 2017
  AND races.name = 'Abu Dhabi Grand Prix'
  AND constructors.constructorId IN (SELECT qualifying.constructorId
									 FROM qualifying
								 		 INNER JOIN races ON races.raceId = qualifying.raceId
									 WHERE races.year = 2017
									   AND races.name = 'Abu Dhabi Grand Prix'
									   AND qualifying.q1 <> ''
                                       AND qualifying.q2 <> ''
                                       AND qualifying.q3 <> ''
									 GROUP BY qualifying.constructorId
									 HAVING COUNT(DISTINCT qualifying.driverId) = 2)
ORDER BY constructors.name ASC
    \end{minted}
    
    \item % Obtener el nombre y apellido de los pilotos y el nombre de aquellas carreras en las que han participado pilotos rusos y polacos.
    \begin{minted}[fontsize=\scriptsize, frame=lines, xleftmargin=\parindent, breaklines, framesep=2mm]{SQL}  
SELECT drivers.forename, drivers.surname, races.name, races.year
FROM drivers
	INNER JOIN results ON results.driverId = drivers.driverId
    INNER JOIN races ON races.raceId = results.raceId
WHERE races.raceId IN (SELECT raceId
					   FROM results 
                           INNER JOIN drivers ON drivers.driverId = results.driverId
					   WHERE drivers.nationality = "Russian")
  AND races.raceId IN (SELECT raceId
					   FROM results 
                           INNER JOIN drivers ON drivers.driverId = results.driverId
					   WHERE drivers.nationality = "Polish")
    \end{minted}
    
    \item % Obtener el nombre y apellidos de los pilotos y el número de vueltas totales recorridas en el año 2011 siempre y cuando sea mayor que la media del número de vueltas totales recorridas el año anterior por todos los pilotos.
    \begin{minted}[fontsize=\scriptsize, frame=lines, xleftmargin=\parindent, breaklines, framesep=2mm]{SQL}  
SELECT drivers.forename, drivers.surname, SUM(results.laps)
FROM drivers
	INNER JOIN results ON results.driverId = drivers.driverId
    INNER JOIN races ON races.raceId = results.raceId
WHERE races.year = 2011
GROUP BY drivers.driverId, drivers.forename, drivers.surname
HAVING SUM(results.laps) > (SELECT AVG(nLaps)
							FROM (SELECT SUM(laps) AS nLaps
								  FROM results
									  INNER JOIN races ON races.raceId = results.raceId
								  WHERE year = 2010
								  GROUP BY driverId) t)
    \end{minted}
    
    \item % Obtener el nombre y año de las carreras en las que se disputó una clasificación (\texttt{qualifying}) pero no se realizaron pitstops.
    \begin{minted}[fontsize=\scriptsize, frame=lines, xleftmargin=\parindent, breaklines, framesep=2mm]{SQL}  
SELECT name, year
FROM races
WHERE raceId IN (SELECT raceId FROM qualifying)
  AND raceId NOT IN (SELECT raceId FROM pitstops)
    \end{minted}
    
    \item % Obtener la nacionalidad de los pilotos que han disputado todas las ediciones del gran premio `Australian Grand Prix'.
    \begin{minted}[fontsize=\scriptsize, frame=lines, xleftmargin=\parindent, breaklines, framesep=2mm]{SQL}  
-- Solución con NOT EXITS
SELECT DISTINCT nationality
FROM drivers
WHERE NOT EXISTS (SELECT *
                  FROM races
                  WHERE name = "Australian Grand Prix"
                    AND NOT EXISTS (SELECT *
                                    FROM results
                                    WHERE results.driverId = drivers.driverId
                                      AND results.raceId = races.raceId))
                                      
-- Solución con COUNT
SELECT drivers.nationality
FROM drivers
	INNER JOIN results ON results.driverId = drivers.driverId
    INNER JOIN races ON races.raceId = results.raceId
WHERE races.name = "Australian Grand Prix"
GROUP BY drivers.driverId, drivers.nationality
HAVING COUNT(*) = (SELECT COUNT(*)
                  FROM races
                  WHERE name = "Australian Grand Prix")
    \end{minted}
       
    \item % Eliminar de la tabla \texttt{qualifying} aquellas tuplas donde un piloto no haya participado en la clasificación.
    \begin{minted}[fontsize=\scriptsize, frame=lines, xleftmargin=\parindent, breaklines, framesep=2mm]{SQL}  
DELETE FROM qualifying
WHERE q1 = '' AND q2 = '' AND q3 = ''
    \end{minted}
            
    \item % Obtener aquellos constructores que habiendo ganado más de 5 carreras entre 2003 y 2010, no hayan participado en ninguna carrera desde el siguiente año.
    \begin{minted}[fontsize=\scriptsize, frame=lines, xleftmargin=\parindent, breaklines, framesep=2mm]{SQL}  
SELECT *
FROM constructors
WHERE constructorId IN (SELECT results.constructorId
						FROM results
							INNER JOIN races ON races.raceId = results.raceId
						WHERE races.year BETWEEN 2003 AND 2010
						  AND results.positionOrder = 1
						GROUP BY results.constructorId
						HAVING COUNT(*) > 5)
  AND constructorId NOT IN (SELECT results.constructorId
							FROM results
								INNER JOIN races ON races.raceId = results.raceId
							WHERE races.year > 2010)
    \end{minted}
        
    \item % Obtener el nombre de la carrera y el año en la que tuvo lugar todos los tipos de incidentes que se enumeran a continuación: descalificación, accidente, colisión, fallo de motor, caja de cambios y transmisión (statusId de 2 a 7).
    \begin{minted}[fontsize=\scriptsize, frame=lines, xleftmargin=\parindent, breaklines, framesep=2mm]{SQL}  
-- Solución con NOT EXISTS
SELECT name, year
FROM races
WHERE NOT EXISTS(SELECT *
                 FROM status
                 WHERE statusId BETWEEN 2 AND 7
                   AND NOT EXISTS(SELECT *
                                  FROM results
                                  WHERE results.raceId = races.raceId
                                    AND results.statusId = status.statusId))
                                    
-- Solución con COUNT
SELECT races.raceId, races.name, races.year, COUNT(DISTINCT results.statusId)
FROM races
	INNER JOIN results ON results.raceId = races.raceId
WHERE results.statusId BETWEEN 2 AND 7
GROUP BY races.raceId, races.name, races.year
HAVING COUNT(DISTINCT results.statusId) = (SELECT COUNT(*)
										   FROM status
                                           WHERE statusId BETWEEN 2 AND 7)
    \end{minted}

    \item % Obtener nombre y apellidos del piloto, el nombre del circuito, el año de la carrera donde un piloto español obtuvo el tiempo de parada más pequeño (atributo \texttt{miliseconds}). Incluya este atributo en la salida de la consulta.
    \begin{minted}[fontsize=\scriptsize, frame=lines, xleftmargin=\parindent, breaklines, framesep=2mm]{SQL}  
SELECT forename, surname, circuits.name, races.year, miliseconds
FROM drivers, pitstops, races, circuits
WHERE drivers.driverId=pitstops.driverId 
  AND pitstops.raceId=races.raceId
  AND races.circuitId=circuits.circuitId
  AND drivers.nationality='Spanish'
  AND miliseconds = (SELECT MIN(miliseconds)
                     FROM pitstops JOIN drivers ON pitstops.driverId=drivers.driverId
                     WHERE drivers.nationality='Spanish');
    \end{minted}
     
    \item % Codifique una consulta que obtenga el nombre de los constructores italianos (\texttt{nationality = `Italian'}) con los que hayan disputado carreras al menos un piloto italiano.
    \begin{minted}[fontsize=\scriptsize, frame=lines, xleftmargin=\parindent, breaklines, framesep=2mm]{SQL}  
SELECT DISTINCT name
FROM constructors
    INNER JOIN results ON results.constructorId = constructors.constructorId
    INNER JOIN drivers ON drivers.driverId = results.resultId
WHERE constructors.nationality LIKE 'Italian'
  AND drivers.nationality LIKE 'Italian'
    \end{minted}
            
    \item % Codifique una consulta que obtenga el nombre y apellidos del piloto que más accidentes (\texttt{status.status = `Accident'}) ha tenido. Mostrar también el número de accidentes.
    \begin{minted}[fontsize=\scriptsize, frame=lines, xleftmargin=\parindent, breaklines, framesep=2mm]{SQL}  
SELECT drivers.forename, drivers.forename, COUNT(*)
FROM drivers
	INNER JOIN results ON results.driverId = drivers.driverId
    INNER JOIN status ON status.statusId = results.statusId
WHERE status.status LIKE 'Accident'
GROUP BY drivers.driverId, drivers.forename, drivers.surname
HAVING COUNT(*) >= ALL (SELECT COUNT(*)
                        FROM results
                        	INNER JOIN status ON status.statusId = results.statusId
                        WHERE status.status LIKE 'Accident'
                        GROUP BY results.driverId) 
    \end{minted}
            
    \item % Codifique una consulta que obtenga el nombre y apellidos de los pilotos que hayan calificado entre los 10 primeros puestos (\texttt{position <= 10}) de todas las carreras del año 2015.
    \begin{minted}[fontsize=\scriptsize, frame=lines, xleftmargin=\parindent, breaklines, framesep=2mm]{SQL}  
SELECT drivers.forename, drivers.surname
FROM drivers
WHERE NOT EXISTS (SELECT *
                  FROM races
                  WHERE year = 2015
                    AND NOT EXISTS (SELECT *
                                    FROM qualifying
                                    WHERE position <= 10
                                      AND qualifying.driverId = drivers.driverId
                                      AND qualifying.raceId = races.raceId))
    \end{minted}
            
    \item % Obtener una lista con los nombres de aquellos constructores italianos (‘Italian’ en inglés) que nunca han competido con pilotos italianos.
    \begin{minted}[fontsize=\scriptsize, frame=lines, xleftmargin=\parindent, breaklines, framesep=2mm]{SQL}  
SELECT name
FROM constructors
WHERE constructorId NOT IN (SELECT constructorId
                            FROM results JOIN drivers ON results.driverId=drivers.driverId
                            WHERE nationality='Italian')
AND nationality='Italian';

    \end{minted}
            
    \item % Obtener toda la información de los constructores que en todas las carreras del año 2006, consiguieron que alguno de sus pilotos quedara entre los diez primeros de la clasificación.
    \begin{minted}[fontsize=\scriptsize, frame=lines, xleftmargin=\parindent, breaklines, framesep=2mm]{SQL}  
SELECT *
FROM constructors C
WHERE NOT EXISTS (SELECT *
                  FROM races R
                  WHERE year=2006
                  AND NOT EXISTS (SELECT *
                                  FROM qualifying
                                  WHERE qualifying.raceId=R.raceId
                                  AND qualifying.constructorId=C.constructorId
                                  AND position <= 10))  

    \end{minted}

    \item % Obtener nombre y apellidos del piloto que acabó la competición con más puntos entre los años 1990 y 2000, así como dicha suma de puntos.
    \begin{minted}[fontsize=\scriptsize, frame=lines, xleftmargin=\parindent, breaklines, framesep=2mm]{SQL}  
SELECT forename, surname, suma
FROM(SELECT driverId, year, SUM(points) as suma
     FROM results JOIN races ON results.raceId=races.raceId
     WHERE year BETWEEN 1990 AND 2000
     GROUP BY driverId, year
     HAVING SUM(points)>0) AS puntuaciones
JOIN drivers ON puntuaciones.driverId=drivers.driverId
WHERE suma >= ALL (SELECT SUM(points)
                   FROM results JOIN races ON results.raceId=races.raceId
                   WHERE year BETWEEN 1990 AND 2000
                   GROUP BY driverId, year
                   HAVING SUM(points)>0)

    \end{minted}
    
    \item % Codifique una consulta que obtenga el nombre y apellidos de los pilotos que ganaron una carrera (\texttt{results.positionOrder = 1}) sin haber estado clasificados entre los 10 primeros pilotos (\texttt{qualifying.position > 10}). Mostrar además el nombre de la carrera y el año en la que lo consiguieron.
    \begin{minted}[fontsize=\scriptsize, frame=lines, xleftmargin=\parindent, breaklines, framesep=2mm]{SQL}  
SELECT DISTINCT d.forename, d.surname, r.name, r.year
FROM drivers d
	INNER JOIN results ON results.driverId = d.driverId
    INNER JOIN races r ON r.raceId = results.raceId
WHERE results.positionOrder = 1
  AND EXISTS (SELECT *
              FROM qualifying
              WHERE qualifying.raceId = r.raceId
                AND qualifying.driverId = d.driverId
                AND position > 10)
    \end{minted}
            
            
    \item % Codifique una consulta que obtenga el nombre y apellidos del piloto que realizó más \emph{pitstops} en una carrera del año 2013. Mostrar también el número de \emph{pitstops}.
    \begin{minted}[fontsize=\scriptsize, frame=lines, xleftmargin=\parindent, breaklines, framesep=2mm]{SQL}  
SELECT drivers.forename, drivers.surname, COUNT(*)
FROM drivers
	INNER JOIN pitstops ON pitstops.driverId = drivers.driverId
    INNER JOIN races ON races.raceId = pitstops.raceId
WHERE races.year = 2013
GROUP BY pitstops.driverId, pitstops.raceId, drivers.forename, drivers.surname
HAVING COUNT(*) >= ALL (SELECT COUNT(*)
                        FROM pitstops
                       	GROUP BY pitstops.driverId, pitstops.raceId)  
    \end{minted}
            
            
    \item % Codifique una consulta que obtenga el nombre y apellidos de los pilotos que hayan quedado entre los 10 primeros puestos (\texttt{positionOrder <= 10}) de todas las carreras del año 2017.
    \begin{minted}[fontsize=\scriptsize, frame=lines, xleftmargin=\parindent, breaklines, framesep=2mm]{SQL}  
SELECT drivers.forename, drivers.surname
FROM drivers
WHERE NOT EXISTS (SELECT *
                  FROM races
                  WHERE year = 2017
                    AND NOT EXISTS (SELECT *
                                    FROM results
                                    WHERE positionOrder <= 10
                                      AND results.driverId = drivers.driverId
                                      AND results.raceId = races.raceId))
    \end{minted}
    
    \item % Obtener toda la información de los constructores que en todas las carreras del año 2006, consiguieron que alguno de sus pilotos quedara entre los diez primeros de la clasificación.
    \begin{minted}[fontsize=\scriptsize, frame=lines, xleftmargin=\parindent, breaklines, framesep=2mm]{SQL}  
SELECT *
FROM constructors C
WHERE NOT EXISTS (SELECT *
                  FROM races R
                  WHERE year=2006
                  AND NOT EXISTS (SELECT *
                                  FROM qualifying
                                  WHERE qualifying.raceId=R.raceId
                                  AND qualifying.constructorId=C.constructorId
                                  AND position <= 10))  

    \end{minted}

    \item % Obtener nombre y apellidos del piloto que acabó la competición con más puntos entre los años 1990 y 2000, así como dicha suma de puntos.
    \begin{minted}[fontsize=\scriptsize, frame=lines, xleftmargin=\parindent, breaklines, framesep=2mm]{SQL}
SELECT forename, surname, suma
FROM(SELECT driverId, year, SUM(points) as suma
     FROM results JOIN races ON results.raceId=races.raceId
     WHERE year BETWEEN 1990 AND 2000
     GROUP BY driverId, year
     HAVING SUM(points)>0) AS puntuaciones
JOIN drivers ON puntuaciones.driverId=drivers.driverId
WHERE suma >= ALL (SELECT SUM(points)
                   FROM results JOIN races ON results.raceId=races.raceId
                   WHERE year BETWEEN 1990 AND 2000
                   GROUP BY driverId, year
                   HAVING SUM(points)>0)

    \end{minted}

    \item % Nombre y apellidos del piloto que realizó más pitstops en una carrera del año 2013. Mostrar también el número de pitstops.
    \begin{minted}[fontsize=\scriptsize, frame=lines, xleftmargin=\parindent, breaklines, framesep=2mm]{SQL}    
    SELECT drivers.forename, drivers.surname, COUNT(*)
FROM drivers
	INNER JOIN pitstops ON pitstops.driverId = drivers.driverId
    INNER JOIN races ON races.raceId = pitstops.raceId
WHERE races.year = 2013
GROUP BY pitstops.driverId, pitstops.raceId, drivers.forename, drivers.surname
HAVING COUNT(*) >= ALL (SELECT COUNT(*)
                        FROM pitstops
                       	GROUP BY pitstops.driverId, pitstops.raceId)  
    \end{minted}
\end{enumerate}


\vspace{2em}
\hrule
\doclicenseThis

\end{document}
